%\title{Reporte Experimento X}
\documentclass[12pt,letterpaper]{article}     % Tipo de documento y otras especificaciones
\usepackage[utf8]{inputenc}                   % Para escribir tildes y eñes
\usepackage[spanish]{babel}                   % Para que los títulos de figuras, tablas y otros estén en español
\addto\captionsspanish{\renewcommand{\tablename}{Tabla}}					% Cambiar nombre a tablas
\addto\captionsspanish{\renewcommand{\listtablename}{Índice de tablas}}		% Cambiar nombre a lista de tablas
\usepackage{geometry}                         
\geometry{left=18mm,right=18mm,top=21mm,bottom=21mm} % Tamaño del área de escritura de la página
\usepackage{ucs}
\usepackage{amsmath}      % Los paquetes ams son desarrollados por la American Mathematical Society
\usepackage{amsfonts}     % y mejoran la escritura de fórmulas y símbolos matemáticos.
\usepackage{amssymb}
\usepackage{graphicx}     % Para insertar gráficas
\usepackage[lofdepth,lotdepth]{subfig}	% Para colocar varias figuras
\usepackage{unitsdef}	  % Para la presentación correcta de unidades
\usepackage{pdfpages}   %incluir paginas de pdf externo, para los anexos
\usepackage{appendix}   %para los anexos
\renewcommand{\unitvaluesep}{\hspace*{4pt}}	% Redimensionamiento del espacio entre magnitud y unidad
\usepackage[colorlinks=true,urlcolor=blue,linkcolor=black,citecolor=black]{hyperref}     % Para insertar hipervínculos y marcadores
\usepackage{float}		% Para ubicar las tablas y figuras justo después del texto
\usepackage{booktabs}	% Para hacer tablas más estilizadas
\batchmode
\bibliographystyle{plain} 
\pagestyle{plain} 
\pagenumbering{arabic}
\usepackage{lastpage}
\usepackage{fancyhdr}	% Para manejar los encabezados y pies de página
\pagestyle{fancy}		% Contenido de los encabezados y pies de pagina
\usepackage{multicol}   % Para varias columnas

%---------------------------Definición del environment resumen---------------------------
\newcounter{resumen}
\setcounter{resumen}{0}
\def\theejemplo{\thechapter.\arabic{resumen}}

\newenvironment{resumen}
{	
	\begin{center}
	\begin{minipage}[t]{500 pt}
	\vspace{5mm}
	\emph{\textbf{Resumen}}
	\\[-2mm]
	\line(1,0){500}
	\\[-4.25 mm]
	\line(1,0){500}
	\\
}
{
	\normalsize
	\\[2mm]
	\footnotesize\textbf{Palabras clave: \footnotesize\@palabras}
	\\[-2mm]
	\line(1,0){500}
	\\[0.5cm]
	\end{minipage}
	\end{center}
}

% -------------------- Para las palabras clave -------------- %
\def\palabras#1{\gdef\@palabras{#1}}

%%%%%%%%%%%%%%%%%%%%%%%%%%%%%%%%%%%%

\lhead{IE0308 - Laboratorio Eléctrico I}
\chead{}
\rhead{Experimento X}	% Aquí va el numero de experimento, al igual que en el titulo
\lfoot{Escuela de Ingeniería Eléctrica}
\cfoot{\thepage\ de \pageref{LastPage}}
\rfoot{Universidad de Costa Rica}


\author{Nombre1, Carné1 \\ Nombre2, Carné2 \\ {\small Grupo 02}\\ Profesor:  \vspace*{3.0in}}
\title{Universidad de Costa Rica\\{\small Facultad de Ingeniería\\Escuela de Ingeniería Eléctrica\\IE0308 – Laboratorio Eléctrico I\\II ciclo 2013\\\vspace*{0.55in} Reporte}\\ Experimento X\vspace*{1.35in}}
%\date{}  				

%%%%%%%%%%%%%%%%% PALABRAS CLAVE 
\palabras{LateX, TINA, Filtros Pasivos}
%%%%%%%%%%%%%%%%%
% Se escriben después del resumen y sintetizan los conceptos fundamentales del experimento a modo de etiquetas


%%%%%%%%%%%%%%%%
\begin{document}	% Inicio del documento
%%%%%%%%%%%%%%%%

\pdfbookmark[1]{Portada}{portada} 	% Marcador para el título

\maketitle							% Título


\tableofcontents
\newpage
\listoffigures
\listoftables

\newpage
\section{Resumen}
%%%%%%%%%%%%%%%%%%%
\begin{resumen}
En esta sección se describe el experimento y la forma en que se desarrollo de manera concis, haciendo enfasis en los resultados y conclusiones más importante en un máximo de 160 palabras.
\end{resumen}
%%%%%%%%%%%%%%%%%%%


%%%%%%%%%%%%%%%%%%%
\section{Objetivos}
%%%%%%%%%%%%%%%%%%%

\subsection{Objetivo general}
%%%%%%%%%%%%%%%%%%%%%%%%%%%%%

Objetivo descrito en la Guía del Experimento.


\subsection{Objetivos específicos}
%%%%%%%%%%%%%%%%%%%%%%%%%%%%%%%%%%

\begin{itemize}
\item Objetivos específicos propuestos por los estudiantes.
\item Son exhaustivos y basados en la nota teórica, el procedimiento y los resultados esperados del experimento.
\end{itemize}

\newpage

%%%%%%%%%%%%%%%%%%%%%%%%%%
\section{Lista de equipos}
%%%%%%%%%%%%%%%%%%%%%%%%%%

La lista de equipos utilizados en el experimento se muestra en la Tabla \ref{T:equipos}.

\begin{table*}[!ht]
\caption{Lista de equipos}
\label{T:equipos}
\begin{center}
\begin{tabular}{| c | c | c | c | c | c | c |}
\hline
\textbf{Equipo} & \multicolumn{3}{| c |}{\textbf{Sesión 1}} & \multicolumn{3}{| c |}{\textbf{Sesión 2}} \\ 
\cline{2-7}
& \textbf{Marca} & \textbf{Modelo} & \textbf{Placa} & \textbf{Marca} & \textbf{Modelo} & \textbf{Placa} \\
\hline
Fuente DC &  &  &  & & &  \\
Osciloscopio digital &  &  &  & & &  \\
Generador de señales &  &  &  & & &  \\
Medidor Multifunción &  &  &  & & &  \\
Protoboard &  &  &  & & &  \\
\hline
\end{tabular}
\end{center}
\end{table*}

\newpage

%%%%%%%%%%%%%%%%%%%%%%%%%%%%%%
\section{Lista de componentes}
%%%%%%%%%%%%%%%%%%%%%%%%%%%%%%

La lista de componentes utilizados en el experimento se muestra en la Tabla \ref{T:componentes}.

\begin{table*} [!ht]
\caption{Lista de componentes}
\label{T:componentes}
\begin{center}
\begin{tabular}{| c | c | c | c | c | c |}
\hline
\textbf{Componente} & \textbf{Sigla} & \textbf{Valor nominal} & \textbf{Valor medido} & \textbf{Tolerancia} & \textbf{Potencia}\\ 
\hline
Resistor cerámico & R1 & \kiloohm[1] & &  & 0,25 W\\ 
Potenciómetro & Rp & \kiloohm[5] & &  & 2 W\\
Capacitor electrolítico & C11 & 150\microfarad & & - & - \\ 
Capacitor cerámico & C & 100\nanofarad & & - & - \\
Diodo & 1N4148 & & & & 500 mW\\
Transistor BJT & 2N2222A & & & & 1000 mW\\
\hline
\end{tabular}
\end{center}
\end{table*}


\newpage
%%%%%%%%%%%%%%%%%%%%%%%
\section{Resultados experimentales y análisis de resultados}
%%%%%%%%%%%%%%%%%%%%%%%

\newpage
%%%%%%%%%%%%%%%%%%%%%%%
\section{Conclusiones y recomendaciones}
%%%%%%%%%%%%%%%%%%%%%%%


%%%%%%%%%%%%%%
% Bibliografía
%%%%%%%%%%%%%%

\newpage

El formato recomendado para la bibliografía es el APA. El siguiente es un ejemplo:

\begin{thebibliography}{}
\bibitem[Apuntes, 2008]{apuntes} \emph{http://www.apuntesdeelectronica.com/componentes/transistor-igbt.htm} consultado el 12/05/2013.
\bibitem[Boylestad, 1998]{Boylestad} Robert L. Boylestad, Louis Nashelsky (1998). \emph{Electronic Devices and Circuit Theory}. New Jersey:  Pearson Prentice Hall, 7th Edition.



\end{thebibliography}


\newpage

%%%%%%%%%%%%%%
% Anexos
%%%%%%%%%%%%%%

\appendix
\section{Anexos}

Incluye información de las hojas del fabricante de los componentes a utilizar o cualquier otra información que se considere necesaria.

La siguiente es una forma de como adjuntar páginas de PDFs con las hojas de fabricante, pero hay muchas formas de como adjuntar los anexos que se consideren necesarios.

%%%%% Transistor NTE123A y NTE159M
\subsection{Transistor 2N2222A}
\includepdf[pages={1-3}]{./Anexos/2N2222A.pdf}



%%%%%%%%%%%%%%
\end{document}
%%%%%%%%%%%%%%